% Theory-owned addition to Supplement S2, 2026-10-30.
% Place AFTER paper/X_symmetry_moved_to_suppl.tex, without editing that source.
% The source includes open editorial \pending markers: these are not resolved here.
% Existing theorem labels and source-specific macros are deliberately reused.
\subsection{Proofs, precise hypotheses, and implementation scope}
\label{sec:suppl-symmetry-proofs-audit}

\paragraph{Exact Fourier orbit and posterior ratio (classical).}
For an image with real channels padded centrally from $256\times256$
to a $384\times384$ periodic canvas, reversing both pixel coordinates
about the center of each original channel maps the padded image $f$
to $f'_{j,k}=f_{-j-1,-k-1}$ (indices taken modulo $384$).
For the two-dimensional discrete Fourier transform, a change of
summation indices gives
\[
 \widehat f'(\omega)=
 e^{2\pi i(\omega_1+\omega_2)/384}\widehat f(-\omega).
\]
Since $f$ is real, $\widehat f(-\omega)=\overline{\widehat f(\omega)}$;
hence $|\widehat f'|=|\widehat f|$. The three channel reversals
commute, forming the eight-element group
$\mathcal G\cong(\mathbb Z/2\mathbb Z)^3$; this is one family of
measurement symmetries, not a uniqueness theorem for the inverse
problem. The derivation assumes the stated even, centrally padded,
periodic operator; asymmetric cropping or another forward operator
must be checked anew.

Write $p(x\mid y)=Z_y^{-1}\pi(x)\ell(y;\mathcal A(x))$ with
$0<Z_y<\infty$. The orbit relation
$\mathcal A(gx)=\mathcal A(x)$ gives the cross-multiplied identity
\[
 \pi(x)p(gx\mid y)=\pi(gx)p(x\mid y),
\]
without dividing by a potentially vanishing density. The ratio in
\cref{prop:ratio} is valid only where the denominators are positive.
Because the group transformations permute coordinates, they preserve
Lebesgue measure. Moreover, conditional on the unordered orbit,
observing $y$ supplies no further orientation information:
if $O=\mathcal GX$, then $Y\perp X\mid O$, because its likelihood is
a function of $O$ alone. This is a standard conditional-independence
consequence, not a new information-theoretic theorem.

\paragraph{Why measurement-only selection is blind.}
For a measurement-based selector with score
$R(x;y)=\varphi(\mathcal A(x),y)$, one has
$R(gx;y)=R(x;y)$. Replacing an individual candidate by its
same-orbit representative therefore cannot change its score or
the selected \emph{index}, provided the tie rule is also invariant.
This extends to likelihood-only particle potentials and pruning,
not to full twisted-SMC proposal/prior importance weights. It also
does not assert that separately generated orbit-related candidates
have identical residuals: their additional reconstruction errors
can be different.

\paragraph{Equivariant transitions and the eight orientations.}
Let an initial law $\mu_T$ be invariant under $\mathcal G$.
If each Markov kernel obeys
$K_i(gx,gB)=K_i(x,B)$, induction shows every later law is invariant.
For a Gaussian kernel $N(m_i(x),s_i^2 I)$, positive
state-independent $s_i$ and $m_i(gx)=g m_i(x)$ suffice. If a
neighborhood $B$ of a target has eight pairwise-disjoint translates
$gB$, their equal probabilities give conditional mass $1/8$ in each,
\emph{provided the orbit has eight distinct elements and the union
has positive probability}. No assertion of native solver equivariance
is implicit in this abstract result.

\paragraph{Proof of the Gaussian-path total-variation inequality.}
Fix $g\in\mathcal G$ and let $P$ be the path law of
$(X_T,\ldots,X_0)$, $Q=g_\#P$. Because $g$ is orthogonal and the
initial law is invariant, the two path laws have the same initial
marginal; their respective conditional transitions are
$N(m_i(x),s_i^2 I)$ and $N(gm_i(g^{-1}x),s_i^2 I)$.
With $\Delta_i^g(x)=m_i(x)-gm_i(g^{-1}x)$, the KL chain rule and
same-covariance Gaussian formula give
\[
 D_{\mathrm{KL}}(P\Vert Q)
 =\sum_i\frac{\mathbb E_P\|\Delta_i^g(X_i)\|_2^2}{2s_i^2}.
\]
Data processing for total variation followed by Pinsker gives
\[
 \operatorname{TV}(L,g_\#L)
 \le\operatorname{TV}(P,Q)
 \le\sqrt{\frac{1}{2}D_{\mathrm{KL}}(P\Vert Q)}
 =\left(\sum_i\frac{\mathbb E_P\|\Delta_i^g(X_i)\|_2^2}
 {4s_i^2}\right)^{1/2}.
\]
The bound is on the law of the state at the end of \emph{this}
Gaussian chain. It assumes the correct, full Markov factorization;
a single Gaussian substep inside a non-Gaussian native update does
not certify the entire transition. Exceptional NaN branches and
stateful optimizers require explicit treatment. A final deterministic
readout $H$ preserves this comparison only if it commutes with $g$;
there is no such general guarantee for a learned terminal denoiser.

For completeness, an arbitrary deterministic readout can violate
it maximally: let $Z\sim N(0,I_2)$ be invariant under coordinate
swap $g$, but apply
$H(z)=(\max(z_1,z_2),\min(z_1,z_2))$.
The output and its swap have disjoint supports (up to a null set),
so their TV distance is one although the preceding Gaussian chain
has zero equivariance defect. Therefore the statement must not be
extended to FAST-DIPS's final $\widehat x_0$ readout by omission.

\paragraph{Posterior--sampler mismatch and a necessary defect.}
For the target posterior $\Pi$ and solver output law $L$,
TV is invariant under measurable bijections. A triangle inequality
therefore yields the classical relation
\[
 \operatorname{TV}(\Pi,g_\#\Pi)
 \le 2\operatorname{TV}(\Pi,L)
       +\operatorname{TV}(L,g_\#L).
\]
For an exactly invariant $L$, a noninvariant target posterior
necessarily incurs at least half its orbit-TV asymmetry in posterior
sampling error. Conversely, for the Gaussian chain above, an output
asymmetry $\eta=|L(B)-L(g^{-1}B)|$ requires
$\sum_i\mathbb E\|\Delta_i^g(X_i)\|^2/s_i^2\ge4\eta^2$.
These are general inequalities, not estimates of the deployed
solvers' defects. Denoiser defects measured on noised clean images
cannot replace the \emph{on-trajectory mean-map defects} here.

\paragraph{Pinned native source audit and scope (C--T3).}
\begin{enumerate}
\item \textbf{DAPS, pinned $\texttt{e7a77d0}$.}
Its data-gradient Langevin updates and isotropic innovations are
structurally equivariant if the denoiser and gradient plumbing are;
the MCMC path also contains deterministic reverse-diffusion segments
and a NaN fallback. The final noise scale in the pinned source is zero,
so the returned image is the final Langevin iterate, not a
new nonzero-noise perturbation. A path-KL proof must account for every
nondegenerate Langevin microstep and intervening deterministic map.
\item \textbf{SITCOM-ODE, pinned $\texttt{275ab67}$.}
The audited phase-retrieval implementation uses ODE score steps,
an SSE-gradient correction, inner optimization and Langevin noise;
it is not interchangeable with the lightweight SITCOM
reimplementation in our code base. Its terminal scheduled
noise scale is likewise zero and it returns the last inner
Langevin iterate. Equivariance requires transformed optimizer states,
initializations, tie rules, and denoiser-Jacobian interactions.
\item \textbf{FAST-DIPS, pinned $\texttt{a9758b38}$.}
The measurement-side ADMM projection and JVP/VJP correction are
structurally equivariant under the pixel-permutation action given
an equivariant denoiser and compatible scalar tolerances. But the
last pixel step returns $\operatorname{clip}(\widehat x_0,-1,1)$,
not the computed final corrected state $x_{0y}$. Its terminal
denoising/clamping is deterministic; the Gaussian-chain TV theorem
can bound a qualifying preterminal law, but \emph{not automatically}
the actual returned reconstruction.
\end{enumerate}
The audit is source-specific, not a proof that any native solver
samples its proposed posterior exactly. It also cannot attribute
all non-equivariance uniquely to the denoiser: exceptional branches,
initialization, and numerical tie handling can matter.

\paragraph{Why circular blur and translation do not remove twins.}
For fixed shift-invariant circular convolution with the \emph{same}
known kernel $k$ for every candidate,
\[
 \mathcal F(x*k)=(\mathcal F x)(\mathcal F k),\qquad
 |\mathcal F((gx)*k)|=|\mathcal Fx|\,|\mathcal Fk|
 =|\mathcal F(x*k)|
\]
whenever $|\mathcal F(gx)|=|\mathcal F x|$. This remains true for
asymmetric $k$: the known blur factor multiplies both twins equally.
A translation contributes only a unit-modulus Fourier phase; unknown
misalignment of that kind cannot identify the orientation either.
This does \emph{not} cover spatially varying blur, boundary cropping,
or a nonconvolutional kernel. Known premeasurement coded illumination
$|\mathcal F(m\odot x)|$ or known reference objects can break the
symmetry, because the mask/reference does not transform with $x$;
these are different measurement operators with separate
identifiability conditions.

% \paragraph{Remaining source/figure checks for paper integration.}
% The authoritative relocated symmetry source still contains pending
% editorial markers for XP--C2 equivariance-defect profiles, an
% ImageNet orbit-error count, a generic-detector result, and the
% exact FFHQ prompt quotation. These are \emph{not inferred or filled}
% from the mathematical arguments above. The planner should resolve
% them from registered empirical outputs or delete the dependent
% sentences before camera-ready integration. All proof statements above
% are valid independently of those pending measurements.
% 